%!TEX root = ../../main.tex
%% ===============================================================
%% 2.7 학습 목적함수의 변형과 정규화
%% 파일: chapters/02_background/07_objectives.tex
%% 상위: chapters/02_background/chapter.tex
%% 정의 라벨: sec:bg-objectives
%% 참조 라벨: -
%% 포함 객체: -
%% 인용: caruana1997multitask, lin2017focal, loshchilov2019decoupled, micikevicius2018mixed, muller2019when, ruder2017overview, szegedy2016rethinking
%% 상태: 초안 (docs/OUTLINE.md 참고)
%% ===============================================================

\section{학습 목적함수의 변형과 정규화}
\label{sec:bg-objectives}

초점 손실(focal loss) \cite{lin2017focal}은 쉬운 표본의 손실 기여를 $(1-p_t)^\gamma$로 줄여 어려운 표본에 학습을 집중시킨다. 라벨 평활화 \cite{szegedy2016rethinking}는 경성 라벨을 균등 분포 쪽으로 $\epsilon$만큼 이동시키며, 그 정규화 효과는 Müller 등 \cite{muller2019when}이 분석하였다. 다중 과제 학습은 Caruana \cite{caruana1997multitask}가 정립하였고 심층 신경망에서의 개관은 Ruder \cite{ruder2017overview}를 따른다. 최적화기는 가중치 감쇠를 기울기 갱신과 분리한 AdamW \cite{loshchilov2019decoupled}를 사용하고, 혼합 정밀도 학습 \cite{micikevicius2018mixed}을 적용한다.
